\chapter{Introducción}

La resolución de problemas es una actividad inherente al ser humano. Desde las decisiones más sencillas de la vida cotidiana hasta el desarrollo de complejos sistemas tecnológicos, las personas necesitan identificar dificultades, analizar alternativas y establecer procedimientos que permitan alcanzar determinados objetivos.

La incorporación de las computadoras en diferentes ámbitos de la sociedad ha ampliado considerablemente las posibilidades de resolver problemas. Sin embargo, una computadora no determina por sí misma qué problema debe solucionar ni qué resultado se considera satisfactorio. Es necesario formular el problema, establecer sus condiciones y desarrollar procedimientos que puedan ser representados y ejecutados computacionalmente.

\citet{farrell2013} aborda la importancia de desarrollar la lógica de los programas antes de proceder a su implementación mediante un lenguaje de programación. Desde esta perspectiva, el diseño de una solución constituye una etapa fundamental para obtener programas organizados y funcionales.

Por ejemplo, si una institución educativa necesita determinar la condición académica de sus estudiantes a partir de sus calificaciones, el proceso no comienza necesariamente escribiendo código. Primero es necesario identificar qué información se requiere, cómo se calcularán los resultados, cuáles son las condiciones de aprobación y qué información deberá proporcionar el sistema.

Este procedimiento aplica estrategias del pensamiento computacional como la descomposición, el reconocimiento de patrones, la abstracción, la generalización y el diseño de algoritmos. Estas habilidades sirven en informática, educación, ingeniería, administración y medicina. La \RefFigura{fig:proceso-resolucion} muestra las etapas del proceso.

\FiguraLibro{01-proceso-resolucion}{Proceso general de resolución de problemas computacionales}{fig:proceso-resolucion}{Elaboración propia.}

El libro utiliza la programación para construir y evaluar soluciones. Los ejemplos emplean Python por su sintaxis accesible. Los principios pueden aplicarse en otros lenguajes.
